% --- Archon physics formalization source begin ---
% archon:physics
% archon:covers PhyXMiniProblems/problem_phyx_mini_0496.lean
% archon:source-report reports/phyx_mini/problem_phyx_mini_0496.source.json

\chapter{Physics problem phyx\_mini\_0496}
\label{ch:PhyXMiniProblems_problem_phyx_mini_0496}

\paragraph{Problem source.}
\#\# Physical scenario

The two 5.0-cm-long parallel electrodes in figure are spaced 1.0 cm apart. A proton enters the plates from one end, an equal distance from both electrodes. A potential difference \textbackslash{}( \textbackslash{}Delta V = 500 \textbackslash{}text\{ V\} \textbackslash{}) across the electrodes deflects the proton so that it strikes the outer end of the lower electrode.

\#\# Question

What magnetic field strength will allow the proton to pass through undeflected while the 500 V potential difference is applied? Assume that both the electric and magnetic fields are confined to the space between the electrodes.

\#\# Answer choices

- A: A: \textbackslash{}( 49 \textbackslash{}text\{ mT\} \textbackslash{})

- B: B: \textbackslash{}( 23 \textbackslash{}text\{ mT\} \textbackslash{})

- C: C: \textbackslash{}( 46 \textbackslash{}text\{ mT\} \textbackslash{})

- D: D: \textbackslash{}( 56 \textbackslash{}text\{ mT\} \textbackslash{})

\#\# Figure caption (auxiliary; use the image as primary evidence)

The image shows a diagram of a charged particle moving between two parallel plates. 

- A positive charge, indicated by a red circle with a plus sign, starts on the left side of the plates.
- The green arrow points horizontally to the right, showing the initial direction of the particle's motion.
- The particle's path curves downward as a blue arrow, indicating the influence of an electric field between the plates.
- Two parallel black lines represent the plates, with a vertical separation labeled as "1.0 cm."
- The horizontal length of the plates is labeled as "5.0 cm."
- Below the diagram, text reads "Trajectory at ΔV = 500 V," suggesting a potential difference between the plates.

\#\# Dataset metadata

Quantum Phenomena; Physical Model Grounding Reasoning, Multi-Formula Reasoning

\paragraph{Recorded answer/context.}
C: C: \textbackslash{}( 46 \textbackslash{}text\{ mT\} \textbackslash{})

\paragraph{Figure/image path.}
/root/proposal\_for\_physic/science-mango/phyx\_mini\_run/phyx\_data/test\_image/496.png

\paragraph{Formalization target.}
create a compiling Lean file with sorry bodies at `PhyXMiniProblems/problem\_phyx\_mini\_0496.lean`. The Lean declarations must preserve the physical quantities, dimensions or dimensional roles, figure labels, governing-law hypotheses, and final relation expressed by this problem.
Use Mathlib/Physlib names found through LeanExplore where available. If a physics API is missing, introduce faithful local abstractions rather than scalar placeholder aliases.

\begin{theorem}[Physics formalization target]
\label{thm:physics:phyx_mini_0496:target}
\lean{PhyXMiniProblems.ProblemPhyXMini0496.problem_phyx_mini_0496}
\uses{def:physics:phyx-mini-0496:phyxminiproblems-problemphyxmini0496-parallelplateprotonsetup, def:physics:phyx-mini-0496:phyxminiproblems-problemphyxmini0496-matchesproblemandfigurereadouts, def:physics:phyx-mini-0496:phyxminiproblems-problemphyxmini0496-hasstandardprotoncalibration, def:physics:phyx-mini-0496:phyxminiproblems-problemphyxmini0496-hasphysicalparameters, def:physics:phyx-mini-0496:phyxminiproblems-problemphyxmini0496-fieldsareconfinedtoplategap, def:physics:phyx-mini-0496:phyxminiproblems-problemphyxmini0496-fieldshaveuniformmagnitudeinplategap, def:physics:phyx-mini-0496:phyxminiproblems-problemphyxmini0496-satisfiesparallelplateelectricfieldlaw, def:physics:phyx-mini-0496:phyxminiproblems-problemphyxmini0496-satisfieselectriconlydeflectionkinematics, def:physics:phyx-mini-0496:phyxminiproblems-problemphyxmini0496-satisfiesundeflectedlorentzforcebalance, def:physics:phyx-mini-0496:phyxminiproblems-problemphyxmini0496-answerstrengthinmilliteslas, def:physics:phyx-mini-0496:phyxminiproblems-problemphyxmini0496-roundstonearestwholemillitesla, def:physics:phyx-mini-0496:phyxminiproblems-problemphyxmini0496-isclosestdisplayedchoice}
The assigned autoformalize agent should translate this physics problem into Lean declarations in the covered file, with theorem and lemma proof bodies written as `by sorry`.
\end{theorem}
\begin{proof}
This is an autoformalization task, not a proof task. Produce statements that can later be proved by the physics prover without weakening the physical model.
\end{proof}

% --- Archon declaration topology begin ---
\section{Declaration topology}
\begin{definition}[Length Quantity]
  \label{def:physics:phyx-mini-0496:phyxminiproblems-problemphyxmini0496-lengthquantity}
  \lean{PhyXMiniProblems.ProblemPhyXMini0496.LengthQuantity}
  A nonnegative physical length
\end{definition}
\begin{proof}
  This is the declared model definition of Length Quantity.
\end{proof}
\begin{definition}[Time Interval Quantity]
  \label{def:physics:phyx-mini-0496:phyxminiproblems-problemphyxmini0496-timeintervalquantity}
  \lean{PhyXMiniProblems.ProblemPhyXMini0496.TimeIntervalQuantity}
  A nonnegative elapsed time
\end{definition}
\begin{proof}
  This is the declared model definition of Time Interval Quantity.
\end{proof}
\begin{definition}[Mass Quantity]
  \label{def:physics:phyx-mini-0496:phyxminiproblems-problemphyxmini0496-massquantity}
  \lean{PhyXMiniProblems.ProblemPhyXMini0496.MassQuantity}
  A nonnegative mass
\end{definition}
\begin{proof}
  This is the declared model definition of Mass Quantity.
\end{proof}
\begin{definition}[Charge Magnitude Quantity]
  \label{def:physics:phyx-mini-0496:phyxminiproblems-problemphyxmini0496-chargemagnitudequantity}
  \lean{PhyXMiniProblems.ProblemPhyXMini0496.ChargeMagnitudeQuantity}
  A nonnegative electric-charge magnitude
\end{definition}
\begin{proof}
  This is the declared model definition of Charge Magnitude Quantity.
\end{proof}
\begin{definition}[speed Dimension]
  \label{def:physics:phyx-mini-0496:phyxminiproblems-problemphyxmini0496-speeddimension}
  \lean{PhyXMiniProblems.ProblemPhyXMini0496.speedDimension}
  The physical dimension L T⁻¹ of speed
\end{definition}
\begin{proof}
  This is the declared model definition of speed Dimension.
\end{proof}
\begin{definition}[potential Difference Dimension]
  \label{def:physics:phyx-mini-0496:phyxminiproblems-problemphyxmini0496-potentialdifferencedimension}
  \lean{PhyXMiniProblems.ProblemPhyXMini0496.potentialDifferenceDimension}
  The physical dimension M L² T⁻² C⁻¹ of potential difference
\end{definition}
\begin{proof}
  This is the declared model definition of potential Difference Dimension.
\end{proof}
\begin{definition}[electric Field Strength Dimension]
  \label{def:physics:phyx-mini-0496:phyxminiproblems-problemphyxmini0496-electricfieldstrengthdimension}
  \lean{PhyXMiniProblems.ProblemPhyXMini0496.electricFieldStrengthDimension}
  The physical dimension M L T⁻² C⁻¹ of electric-field strength
\end{definition}
\begin{proof}
  This is the declared model definition of electric Field Strength Dimension.
\end{proof}
\begin{definition}[magnetic Field Strength Dimension]
  \label{def:physics:phyx-mini-0496:phyxminiproblems-problemphyxmini0496-magneticfieldstrengthdimension}
  \lean{PhyXMiniProblems.ProblemPhyXMini0496.magneticFieldStrengthDimension}
  The physical dimension M T⁻¹ C⁻¹ of magnetic-field strength
\end{definition}
\begin{proof}
  This is the declared model definition of magnetic Field Strength Dimension.
\end{proof}
\begin{definition}[Speed Quantity]
  \label{def:physics:phyx-mini-0496:phyxminiproblems-problemphyxmini0496-speedquantity}
  \lean{PhyXMiniProblems.ProblemPhyXMini0496.SpeedQuantity}
  \uses{def:physics:phyx-mini-0496:phyxminiproblems-problemphyxmini0496-speeddimension}
  A nonnegative speed magnitude
\end{definition}
\begin{proof}
  This is the declared model definition of Speed Quantity.
\end{proof}
\begin{definition}[Potential Difference Quantity]
  \label{def:physics:phyx-mini-0496:phyxminiproblems-problemphyxmini0496-potentialdifferencequantity}
  \lean{PhyXMiniProblems.ProblemPhyXMini0496.PotentialDifferenceQuantity}
  \uses{def:physics:phyx-mini-0496:phyxminiproblems-problemphyxmini0496-potentialdifferencedimension}
  A nonnegative potential-difference magnitude
\end{definition}
\begin{proof}
  This is the declared model definition of Potential Difference Quantity.
\end{proof}
\begin{definition}[Electric Field Strength Quantity]
  \label{def:physics:phyx-mini-0496:phyxminiproblems-problemphyxmini0496-electricfieldstrengthquantity}
  \lean{PhyXMiniProblems.ProblemPhyXMini0496.ElectricFieldStrengthQuantity}
  \uses{def:physics:phyx-mini-0496:phyxminiproblems-problemphyxmini0496-electricfieldstrengthdimension}
  A nonnegative electric-field magnitude
\end{definition}
\begin{proof}
  This is the declared model definition of Electric Field Strength Quantity.
\end{proof}
\begin{definition}[Magnetic Field Strength Quantity]
  \label{def:physics:phyx-mini-0496:phyxminiproblems-problemphyxmini0496-magneticfieldstrengthquantity}
  \lean{PhyXMiniProblems.ProblemPhyXMini0496.MagneticFieldStrengthQuantity}
  \uses{def:physics:phyx-mini-0496:phyxminiproblems-problemphyxmini0496-magneticfieldstrengthdimension}
  A nonnegative magnetic-field magnitude
\end{definition}
\begin{proof}
  This is the declared model definition of Magnetic Field Strength Quantity.
\end{proof}
\begin{definition}[length In Meters]
  \label{def:physics:phyx-mini-0496:phyxminiproblems-problemphyxmini0496-lengthinmeters}
  \lean{PhyXMiniProblems.ProblemPhyXMini0496.lengthInMeters}
  \uses{def:physics:phyx-mini-0496:phyxminiproblems-problemphyxmini0496-lengthquantity}
  SI length readout in metres
\end{definition}
\begin{proof}
  This is the declared model definition of length In Meters.
\end{proof}
\begin{definition}[time In Seconds]
  \label{def:physics:phyx-mini-0496:phyxminiproblems-problemphyxmini0496-timeinseconds}
  \lean{PhyXMiniProblems.ProblemPhyXMini0496.timeInSeconds}
  \uses{def:physics:phyx-mini-0496:phyxminiproblems-problemphyxmini0496-timeintervalquantity}
  SI elapsed-time readout in seconds
\end{definition}
\begin{proof}
  This is the declared model definition of time In Seconds.
\end{proof}
\begin{definition}[mass In Kilograms]
  \label{def:physics:phyx-mini-0496:phyxminiproblems-problemphyxmini0496-massinkilograms}
  \lean{PhyXMiniProblems.ProblemPhyXMini0496.massInKilograms}
  \uses{def:physics:phyx-mini-0496:phyxminiproblems-problemphyxmini0496-massquantity}
  SI mass readout in kilograms
\end{definition}
\begin{proof}
  This is the declared model definition of mass In Kilograms.
\end{proof}
\begin{definition}[charge In Coulombs]
  \label{def:physics:phyx-mini-0496:phyxminiproblems-problemphyxmini0496-chargeincoulombs}
  \lean{PhyXMiniProblems.ProblemPhyXMini0496.chargeInCoulombs}
  \uses{def:physics:phyx-mini-0496:phyxminiproblems-problemphyxmini0496-chargemagnitudequantity}
  SI charge-magnitude readout in coulombs
\end{definition}
\begin{proof}
  This is the declared model definition of charge In Coulombs.
\end{proof}
\begin{definition}[speed In Meters Per Second]
  \label{def:physics:phyx-mini-0496:phyxminiproblems-problemphyxmini0496-speedinmeterspersecond}
  \lean{PhyXMiniProblems.ProblemPhyXMini0496.speedInMetersPerSecond}
  \uses{def:physics:phyx-mini-0496:phyxminiproblems-problemphyxmini0496-speedquantity}
  SI speed readout in metres per second
\end{definition}
\begin{proof}
  This is the declared model definition of speed In Meters Per Second.
\end{proof}
\begin{definition}[potential Difference In Volts]
  \label{def:physics:phyx-mini-0496:phyxminiproblems-problemphyxmini0496-potentialdifferenceinvolts}
  \lean{PhyXMiniProblems.ProblemPhyXMini0496.potentialDifferenceInVolts}
  \uses{def:physics:phyx-mini-0496:phyxminiproblems-problemphyxmini0496-potentialdifferencequantity}
  SI potential-difference readout in volts
\end{definition}
\begin{proof}
  This is the declared model definition of potential Difference In Volts.
\end{proof}
\begin{definition}[electric Field Strength In Volts Per Meter]
  \label{def:physics:phyx-mini-0496:phyxminiproblems-problemphyxmini0496-electricfieldstrengthinvoltspermeter}
  \lean{PhyXMiniProblems.ProblemPhyXMini0496.electricFieldStrengthInVoltsPerMeter}
  \uses{def:physics:phyx-mini-0496:phyxminiproblems-problemphyxmini0496-electricfieldstrengthquantity}
  SI electric-field-strength readout in volts per metre
\end{definition}
\begin{proof}
  This is the declared model definition of electric Field Strength In Volts Per Meter.
\end{proof}
\begin{definition}[magnetic Field Strength In Teslas]
  \label{def:physics:phyx-mini-0496:phyxminiproblems-problemphyxmini0496-magneticfieldstrengthinteslas}
  \lean{PhyXMiniProblems.ProblemPhyXMini0496.magneticFieldStrengthInTeslas}
  \uses{def:physics:phyx-mini-0496:phyxminiproblems-problemphyxmini0496-magneticfieldstrengthquantity}
  SI magnetic-field-strength readout in teslas
\end{definition}
\begin{proof}
  This is the declared model definition of magnetic Field Strength In Teslas.
\end{proof}
\begin{definition}[magnetic Field Strength In Milliteslas]
  \label{def:physics:phyx-mini-0496:phyxminiproblems-problemphyxmini0496-magneticfieldstrengthinmilliteslas}
  \lean{PhyXMiniProblems.ProblemPhyXMini0496.magneticFieldStrengthInMilliteslas}
  \uses{def:physics:phyx-mini-0496:phyxminiproblems-problemphyxmini0496-magneticfieldstrengthquantity, def:physics:phyx-mini-0496:phyxminiproblems-problemphyxmini0496-magneticfieldstrengthinteslas}
  Magnetic-field-strength readout in milliteslas
\end{definition}
\begin{proof}
  This is the declared model definition of magnetic Field Strength In Milliteslas.
\end{proof}
\begin{definition}[Particle Species]
  \label{def:physics:phyx-mini-0496:phyxminiproblems-problemphyxmini0496-particlespecies}
  \lean{PhyXMiniProblems.ProblemPhyXMini0496.ParticleSpecies}
  Charged-particle species distinguished by the model
\end{definition}
\begin{proof}
  This is the declared model definition of Particle Species.
\end{proof}
\begin{definition}[Electrode Label]
  \label{def:physics:phyx-mini-0496:phyxminiproblems-problemphyxmini0496-electrodelabel}
  \lean{PhyXMiniProblems.ProblemPhyXMini0496.ElectrodeLabel}
  The upper and lower electrode labels from the supplied figure
\end{definition}
\begin{proof}
  This is the declared model definition of Electrode Label.
\end{proof}
\begin{definition}[Electrode Arrangement]
  \label{def:physics:phyx-mini-0496:phyxminiproblems-problemphyxmini0496-electrodearrangement}
  \lean{PhyXMiniProblems.ProblemPhyXMini0496.ElectrodeArrangement}
  Whether the long electrodes have the parallel geometry shown
\end{definition}
\begin{proof}
  This is the declared model definition of Electrode Arrangement.
\end{proof}
\begin{definition}[Spatial Direction]
  \label{def:physics:phyx-mini-0496:phyxminiproblems-problemphyxmini0496-spatialdirection}
  \lean{PhyXMiniProblems.ProblemPhyXMini0496.SpatialDirection}
  Named directions used in the side-view diagram
\end{definition}
\begin{proof}
  This is the declared model definition of Spatial Direction.
\end{proof}
\begin{definition}[Parallel Plate Proton Setup]
  \label{def:physics:phyx-mini-0496:phyxminiproblems-problemphyxmini0496-parallelplateprotonsetup}
  \lean{PhyXMiniProblems.ProblemPhyXMini0496.ParallelPlateProtonSetup}
  \uses{def:physics:phyx-mini-0496:phyxminiproblems-problemphyxmini0496-lengthquantity, def:physics:phyx-mini-0496:phyxminiproblems-problemphyxmini0496-timeintervalquantity, def:physics:phyx-mini-0496:phyxminiproblems-problemphyxmini0496-massquantity, def:physics:phyx-mini-0496:phyxminiproblems-problemphyxmini0496-chargemagnitudequantity, def:physics:phyx-mini-0496:phyxminiproblems-problemphyxmini0496-speedquantity, def:physics:phyx-mini-0496:phyxminiproblems-problemphyxmini0496-potentialdifferencequantity, def:physics:phyx-mini-0496:phyxminiproblems-problemphyxmini0496-electricfieldstrengthquantity, def:physics:phyx-mini-0496:phyxminiproblems-problemphyxmini0496-magneticfieldstrengthquantity, def:physics:phyx-mini-0496:phyxminiproblems-problemphyxmini0496-particlespecies, def:physics:phyx-mini-0496:phyxminiproblems-problemphyxmini0496-electrodelabel, def:physics:phyx-mini-0496:phyxminiproblems-problemphyxmini0496-electrodearrangement, def:physics:phyx-mini-0496:phyxminiproblems-problemphyxmini0496-spatialdirection}
  Physical quantities and diagrammatic observables for the experiment. The magnetic strength is an independent physical observable. Its requested numerical value is not stored here and is constrained only by the general laws below
\end{definition}
\begin{proof}
  This is the declared model definition of Parallel Plate Proton Setup.
\end{proof}
\begin{definition}[Matches Problem And Figure Readouts]
  \label{def:physics:phyx-mini-0496:phyxminiproblems-problemphyxmini0496-matchesproblemandfigurereadouts}
  \lean{PhyXMiniProblems.ProblemPhyXMini0496.MatchesProblemAndFigureReadouts}
  \uses{def:physics:phyx-mini-0496:phyxminiproblems-problemphyxmini0496-lengthinmeters, def:physics:phyx-mini-0496:phyxminiproblems-problemphyxmini0496-potentialdifferenceinvolts, def:physics:phyx-mini-0496:phyxminiproblems-problemphyxmini0496-parallelplateprotonsetup}
  Numerical and qualitative data supplied by the problem and primary image. The 0.005 m entry height records that the proton starts equally far from the plates. The impact fields record the electric-only trajectory reaching the outer end of the lower electrode. No magnetic answer value occurs here
\end{definition}
\begin{proof}
  This is the declared model definition of Matches Problem And Figure Readouts.
\end{proof}
\begin{definition}[Has Standard Proton Calibration]
  \label{def:physics:phyx-mini-0496:phyxminiproblems-problemphyxmini0496-hasstandardprotoncalibration}
  \lean{PhyXMiniProblems.ProblemPhyXMini0496.HasStandardProtonCalibration}
  \uses{def:physics:phyx-mini-0496:phyxminiproblems-problemphyxmini0496-massinkilograms, def:physics:phyx-mini-0496:phyxminiproblems-problemphyxmini0496-chargeincoulombs, def:physics:phyx-mini-0496:phyxminiproblems-problemphyxmini0496-parallelplateprotonsetup}
  Standard proton mass and elementary-charge readouts used for evaluating the multiple-choice result. These are calibration data, not the requested field
\end{definition}
\begin{proof}
  This is the declared model definition of Has Standard Proton Calibration.
\end{proof}
\begin{definition}[Has Physical Parameters]
  \label{def:physics:phyx-mini-0496:phyxminiproblems-problemphyxmini0496-hasphysicalparameters}
  \lean{PhyXMiniProblems.ProblemPhyXMini0496.HasPhysicalParameters}
  \uses{def:physics:phyx-mini-0496:phyxminiproblems-problemphyxmini0496-lengthinmeters, def:physics:phyx-mini-0496:phyxminiproblems-problemphyxmini0496-timeinseconds, def:physics:phyx-mini-0496:phyxminiproblems-problemphyxmini0496-massinkilograms, def:physics:phyx-mini-0496:phyxminiproblems-problemphyxmini0496-chargeincoulombs, def:physics:phyx-mini-0496:phyxminiproblems-problemphyxmini0496-speedinmeterspersecond, def:physics:phyx-mini-0496:phyxminiproblems-problemphyxmini0496-electricfieldstrengthinvoltspermeter, def:physics:phyx-mini-0496:phyxminiproblems-problemphyxmini0496-magneticfieldstrengthinteslas, def:physics:phyx-mini-0496:phyxminiproblems-problemphyxmini0496-parallelplateprotonsetup}
  Positivity and non-vacuity conditions for the physical experiment
\end{definition}
\begin{proof}
  This is the declared model definition of Has Physical Parameters.
\end{proof}
\begin{definition}[Field Is Confined To]
  \label{def:physics:phyx-mini-0496:phyxminiproblems-problemphyxmini0496-fieldisconfinedto}
  \lean{PhyXMiniProblems.ProblemPhyXMini0496.FieldIsConfinedTo}
  A spacetime-dependent vector field vanishes outside the specified region
\end{definition}
\begin{proof}
  This is the declared model definition of Field Is Confined To.
\end{proof}
\begin{definition}[Fields Are Confined To Plate Gap]
  \label{def:physics:phyx-mini-0496:phyxminiproblems-problemphyxmini0496-fieldsareconfinedtoplategap}
  \lean{PhyXMiniProblems.ProblemPhyXMini0496.FieldsAreConfinedToPlateGap}
  \uses{def:physics:phyx-mini-0496:phyxminiproblems-problemphyxmini0496-parallelplateprotonsetup, def:physics:phyx-mini-0496:phyxminiproblems-problemphyxmini0496-fieldisconfinedto}
  Both applied fields are confined to the gap between the electrodes
\end{definition}
\begin{proof}
  This is the declared model definition of Fields Are Confined To Plate Gap.
\end{proof}
\begin{definition}[Fields Have Uniform Magnitude In Plate Gap]
  \label{def:physics:phyx-mini-0496:phyxminiproblems-problemphyxmini0496-fieldshaveuniformmagnitudeinplategap}
  \lean{PhyXMiniProblems.ProblemPhyXMini0496.FieldsHaveUniformMagnitudeInPlateGap}
  \uses{def:physics:phyx-mini-0496:phyxminiproblems-problemphyxmini0496-electricfieldstrengthinvoltspermeter, def:physics:phyx-mini-0496:phyxminiproblems-problemphyxmini0496-magneticfieldstrengthinteslas, def:physics:phyx-mini-0496:phyxminiproblems-problemphyxmini0496-parallelplateprotonsetup}
  The dimensionful scalar strengths are the uniform magnitudes of the Physlib vector fields throughout the plate gap
\end{definition}
\begin{proof}
  This is the declared model definition of Fields Have Uniform Magnitude In Plate Gap.
\end{proof}
\begin{definition}[Satisfies Parallel Plate Electric Field Law]
  \label{def:physics:phyx-mini-0496:phyxminiproblems-problemphyxmini0496-satisfiesparallelplateelectricfieldlaw}
  \lean{PhyXMiniProblems.ProblemPhyXMini0496.SatisfiesParallelPlateElectricFieldLaw}
  \uses{def:physics:phyx-mini-0496:phyxminiproblems-problemphyxmini0496-lengthinmeters, def:physics:phyx-mini-0496:phyxminiproblems-problemphyxmini0496-potentialdifferenceinvolts, def:physics:phyx-mini-0496:phyxminiproblems-problemphyxmini0496-electricfieldstrengthinvoltspermeter, def:physics:phyx-mini-0496:phyxminiproblems-problemphyxmini0496-parallelplateprotonsetup}
  Uniform parallel-plate electrostatics: E = ΔV / d
\end{definition}
\begin{proof}
  This is the declared model definition of Satisfies Parallel Plate Electric Field Law.
\end{proof}
\begin{definition}[Satisfies Electric Only Deflection Kinematics]
  \label{def:physics:phyx-mini-0496:phyxminiproblems-problemphyxmini0496-satisfieselectriconlydeflectionkinematics}
  \lean{PhyXMiniProblems.ProblemPhyXMini0496.SatisfiesElectricOnlyDeflectionKinematics}
  \uses{def:physics:phyx-mini-0496:phyxminiproblems-problemphyxmini0496-lengthinmeters, def:physics:phyx-mini-0496:phyxminiproblems-problemphyxmini0496-timeinseconds, def:physics:phyx-mini-0496:phyxminiproblems-problemphyxmini0496-massinkilograms, def:physics:phyx-mini-0496:phyxminiproblems-problemphyxmini0496-chargeincoulombs, def:physics:phyx-mini-0496:phyxminiproblems-problemphyxmini0496-speedinmeterspersecond, def:physics:phyx-mini-0496:phyxminiproblems-problemphyxmini0496-electricfieldstrengthinvoltspermeter, def:physics:phyx-mini-0496:phyxminiproblems-problemphyxmini0496-parallelplateprotonsetup}
  Horizontal uniform motion and downward constant-acceleration motion during the electric-only traversal. The vertical equation is y = (1/2) (q E / m) t², with zero initial vertical velocity
\end{definition}
\begin{proof}
  This is the declared model definition of Satisfies Electric Only Deflection Kinematics.
\end{proof}
\begin{definition}[Satisfies Undeflected Lorentz Force Balance]
  \label{def:physics:phyx-mini-0496:phyxminiproblems-problemphyxmini0496-satisfiesundeflectedlorentzforcebalance}
  \lean{PhyXMiniProblems.ProblemPhyXMini0496.SatisfiesUndeflectedLorentzForceBalance}
  \uses{def:physics:phyx-mini-0496:phyxminiproblems-problemphyxmini0496-chargeincoulombs, def:physics:phyx-mini-0496:phyxminiproblems-problemphyxmini0496-speedinmeterspersecond, def:physics:phyx-mini-0496:phyxminiproblems-problemphyxmini0496-electricfieldstrengthinvoltspermeter, def:physics:phyx-mini-0496:phyxminiproblems-problemphyxmini0496-magneticfieldstrengthinteslas, def:physics:phyx-mini-0496:phyxminiproblems-problemphyxmini0496-parallelplateprotonsetup}
  For a right-moving positive proton and a field into the page, the upward magnetic force cancels the downward electric force. The transverse Lorentz force balance is q E = q v B; it contains no numerical answer for B
\end{definition}
\begin{proof}
  This is the declared model definition of Satisfies Undeflected Lorentz Force Balance.
\end{proof}
\begin{definition}[Answer Choice]
  \label{def:physics:phyx-mini-0496:phyxminiproblems-problemphyxmini0496-answerchoice}
  \lean{PhyXMiniProblems.ProblemPhyXMini0496.AnswerChoice}
  The four answer labels printed with the problem
\end{definition}
\begin{proof}
  This is the declared model definition of Answer Choice.
\end{proof}
\begin{definition}[answer Strength In Milliteslas]
  \label{def:physics:phyx-mini-0496:phyxminiproblems-problemphyxmini0496-answerstrengthinmilliteslas}
  \lean{PhyXMiniProblems.ProblemPhyXMini0496.answerStrengthInMilliteslas}
  \uses{def:physics:phyx-mini-0496:phyxminiproblems-problemphyxmini0496-answerchoice}
  Magnetic strength in milliteslas displayed by an answer choice
\end{definition}
\begin{proof}
  This is the declared model definition of answer Strength In Milliteslas.
\end{proof}
\begin{definition}[Rounds To Nearest Whole Millitesla]
  \label{def:physics:phyx-mini-0496:phyxminiproblems-problemphyxmini0496-roundstonearestwholemillitesla}
  \lean{PhyXMiniProblems.ProblemPhyXMini0496.RoundsToNearestWholeMillitesla}
  \uses{def:physics:phyx-mini-0496:phyxminiproblems-problemphyxmini0496-magneticfieldstrengthinmilliteslas, def:physics:phyx-mini-0496:phyxminiproblems-problemphyxmini0496-parallelplateprotonsetup}
  The physical strength is within half a millitesla of a whole-mT readout
\end{definition}
\begin{proof}
  This is the declared model definition of Rounds To Nearest Whole Millitesla.
\end{proof}
\begin{definition}[Is Closest Displayed Choice]
  \label{def:physics:phyx-mini-0496:phyxminiproblems-problemphyxmini0496-isclosestdisplayedchoice}
  \lean{PhyXMiniProblems.ProblemPhyXMini0496.IsClosestDisplayedChoice}
  \uses{def:physics:phyx-mini-0496:phyxminiproblems-problemphyxmini0496-magneticfieldstrengthinmilliteslas, def:physics:phyx-mini-0496:phyxminiproblems-problemphyxmini0496-parallelplateprotonsetup, def:physics:phyx-mini-0496:phyxminiproblems-problemphyxmini0496-answerchoice, def:physics:phyx-mini-0496:phyxminiproblems-problemphyxmini0496-answerstrengthinmilliteslas}
  A displayed answer is at least as close as every other displayed answer
\end{definition}
\begin{proof}
  This is the declared model definition of Is Closest Displayed Choice.
\end{proof}
\begin{lemma}[electric Field Strength eq fifty Thousand volts Per Meter]
  \label{lem:physics:phyx-mini-0496:phyxminiproblems-problemphyxmini0496-electricfieldstrength-eq-fiftythousand-voltspermeter}
  \lean{PhyXMiniProblems.ProblemPhyXMini0496.electricFieldStrength_eq_fiftyThousand_voltsPerMeter}
  \uses{def:physics:phyx-mini-0496:phyxminiproblems-problemphyxmini0496-electricfieldstrengthinvoltspermeter, def:physics:phyx-mini-0496:phyxminiproblems-problemphyxmini0496-parallelplateprotonsetup, def:physics:phyx-mini-0496:phyxminiproblems-problemphyxmini0496-matchesproblemandfigurereadouts, def:physics:phyx-mini-0496:phyxminiproblems-problemphyxmini0496-satisfiesparallelplateelectricfieldlaw}
  The 500 V across 1.0 cm plates gives 50000 V/m
\end{lemma}
\begin{proof}
  The conclusion follows from the governing relations and figure readouts named in the statement of electric Field Strength eq fifty Thousand volts Per Meter.
\end{proof}
\begin{lemma}[required Magnetic Field exact Formula]
  \label{lem:physics:phyx-mini-0496:phyxminiproblems-problemphyxmini0496-requiredmagneticfield-exactformula}
  \lean{PhyXMiniProblems.ProblemPhyXMini0496.requiredMagneticField_exactFormula}
  \uses{def:physics:phyx-mini-0496:phyxminiproblems-problemphyxmini0496-lengthinmeters, def:physics:phyx-mini-0496:phyxminiproblems-problemphyxmini0496-massinkilograms, def:physics:phyx-mini-0496:phyxminiproblems-problemphyxmini0496-chargeincoulombs, def:physics:phyx-mini-0496:phyxminiproblems-problemphyxmini0496-potentialdifferenceinvolts, def:physics:phyx-mini-0496:phyxminiproblems-problemphyxmini0496-magneticfieldstrengthinteslas, def:physics:phyx-mini-0496:phyxminiproblems-problemphyxmini0496-parallelplateprotonsetup, def:physics:phyx-mini-0496:phyxminiproblems-problemphyxmini0496-matchesproblemandfigurereadouts, def:physics:phyx-mini-0496:phyxminiproblems-problemphyxmini0496-hasphysicalparameters, def:physics:phyx-mini-0496:phyxminiproblems-problemphyxmini0496-satisfiesparallelplateelectricfieldlaw, def:physics:phyx-mini-0496:phyxminiproblems-problemphyxmini0496-satisfieselectriconlydeflectionkinematics, def:physics:phyx-mini-0496:phyxminiproblems-problemphyxmini0496-satisfiesundeflectedlorentzforcebalance}
  Eliminating transit time and speed from the electric trajectory and Lorentz balance gives B = sqrt (m ΔV / q) / L
\end{lemma}
\begin{proof}
  The conclusion follows from the governing relations and figure readouts named in the statement of required Magnetic Field exact Formula.
\end{proof}
\begin{lemma}[required Magnetic Field millitesla Bounds]
  \label{lem:physics:phyx-mini-0496:phyxminiproblems-problemphyxmini0496-requiredmagneticfield-milliteslabounds}
  \lean{PhyXMiniProblems.ProblemPhyXMini0496.requiredMagneticField_milliteslaBounds}
  \uses{def:physics:phyx-mini-0496:phyxminiproblems-problemphyxmini0496-magneticfieldstrengthinmilliteslas, def:physics:phyx-mini-0496:phyxminiproblems-problemphyxmini0496-parallelplateprotonsetup, def:physics:phyx-mini-0496:phyxminiproblems-problemphyxmini0496-matchesproblemandfigurereadouts, def:physics:phyx-mini-0496:phyxminiproblems-problemphyxmini0496-hasstandardprotoncalibration, def:physics:phyx-mini-0496:phyxminiproblems-problemphyxmini0496-hasphysicalparameters, def:physics:phyx-mini-0496:phyxminiproblems-problemphyxmini0496-satisfiesparallelplateelectricfieldlaw, def:physics:phyx-mini-0496:phyxminiproblems-problemphyxmini0496-satisfieselectriconlydeflectionkinematics, def:physics:phyx-mini-0496:phyxminiproblems-problemphyxmini0496-satisfiesundeflectedlorentzforcebalance}
  The calibrated exact result lies between 45.5 mT and 46.5 mT
\end{lemma}
\begin{proof}
  The conclusion follows from the governing relations and figure readouts named in the statement of required Magnetic Field millitesla Bounds.
\end{proof}
% --- Archon declaration topology end ---
% --- Archon physics formalization source end ---
